\documentclass[11pt]{article}
\usepackage[margin=1in]{geometry}
\usepackage{amsmath,amssymb,amsthm,mathtools,booktabs,hyperref}
\hypersetup{colorlinks=true,linkcolor=blue,urlcolor=blue,citecolor=blue}
\newtheorem{theorem}{Theorem}
\newtheorem{lemma}[theorem]{Lemma}
\newtheorem{corollary}[theorem]{Corollary}
\newtheorem{proposition}[theorem]{Proposition}
\theoremstyle{remark}\newtheorem{remark}[theorem]{Remark}
\newcommand{\sgn}{\operatorname{sign}}
\newcommand{\E}{\mathbb E}\newcommand{\Pp}{\mathbb P}
\title{Exponentially many online mistakes for an isotropic\protect\\ strong-prior Bayesian tanh network}
\author{Dmitry Vaintrob \and Kaarel H\"anni \and ChatGPT-6 Astra}
\date{November 2026}
\hypersetup{colorlinks=true,linkcolor=blue,urlcolor=blue,citecolor=blue,
 pdftitle={Exponentially many online mistakes for an isotropic strong-prior Bayesian tanh network},
 pdfauthor={Dmitry Vaintrob; Kaarel Hanni; ChatGPT-6 Astra}}
\begin{document}\maketitle
\begin{abstract}
We give a fixed input order on which exact Bayesian prediction with an
isotropic Gaussian tanh-network prior makes exponentially many mistakes
in the strong-prior limit. The prior is conditioned on an absolute output
margin at every Boolean input and then on exact training labels. The learner
has one neuron per hidden layer. On $d=2k+1$ input bits, a comparator with a
selector bit forces $2^k-k-1$ mistakes after only $O(k)$ calibration examples.
All these mistakes occur with arbitrarily high probability at every
sufficiently small positive prior scale.
The result covers posterior sampling, posterior-label majority and a
strict $2/3$ confidence criterion. The target has a $\Theta(k)$-gate Boolean
circuit and, at depth $2k+2$, a width-one tanh implementation with weights
bounded by two. Under the convention that teacher complexity $C_L$ bounds
both width and individual weights, this implementation has $C_L=2$ while
the learner width is $N=1$. Thus the result does not settle the separate
question requiring $N\ge C_L$, or a fixed-depth counterexample with
polynomial bounded-weight teacher complexity. A separate appendix proves
an all-width analogue when the activation is the identity.
\end{abstract}

\section{Model and result}

A consistent Bayesian learner need not make few mistakes when inputs arrive
in an unfavorable order. We establish this phenomenon for an unchanged
isotropic Gaussian neural prior, using its full posterior rather than a
choice among maximum-probability classifiers.

The input domain is $\{-1,1\}^d$. There is a constant input coordinate,
no separate internal biases, one neuron in each of $L-1$ hidden layers,
and a linear readout. For $L\ge2$, define the score recursively by
\[
 \begin{aligned}
 h_1(x)&=\tanh\!\left(a_0+\sum_{j=1}^d a_jx_j\right),\\
 h_{j+1}(x)&=\tanh(r_jh_j(x))\quad(1\le j\le L-2),\\
 z_\theta(x)&=b h_{L-1}(x).
\end{aligned}
\]
Thus the list of internal scalar weights is empty when $L=2$.
Independently,
\[
 a_j\sim N(0,\sigma^2/d),\qquad r_j,b\sim N(0,\sigma^2).
\]
The prior is conditioned on $|z_\theta(x)|>1/2$ for every input.
After each prediction, the true label is revealed, and the posterior is
conditioned on its exact sign. The strong-prior limit sends $\sigma$ to zero
at fixed $d,L$ and width. The variances of the first-layer coordinates are
exactly equal; no anisotropic perturbation is used.

Let $p_t$ be the posterior probability of the correct label before round
$t$. Posterior sampling makes a mistake with probability $1-p_t$.
Posterior-label majority is wrong when $p_t<1/2$, and the strict $2/3$
confidence criterion fails when $p_t\le2/3$.

\begin{theorem}[Exponential online mistakes]\label{thm:main}
Fix $k\ge3$, put $d=2k+1$, and take any fixed scalar depth $L\ge2$.
Write an input as $(x,y,z)$ with $x,y\in\{-1,1\}^k$ and $z\in\{-1,1\}$,
and set
\[
 X=\sum_{i=0}^{k-1}2^i\frac{x_i+1}{2},\qquad
 Y=\sum_{i=0}^{k-1}2^i\frac{y_i+1}{2}.
\]
Consider the target
\[
 f(x,y,z)=+1\quad\Longleftrightarrow\quad z=-1\text{ and }X\ge Y.
\]
There is a fixed permutation of the input cube, beginning with at most
$4k+3$ calibration inputs, and a specified set of
\[
 T_k=2^k-k-1\ \ge\ 2^{k-1}
\]
subsequent rounds on which the correct-label posterior probability tends
to zero as $\sigma\downarrow0$.

Consequently, for every $\eta\in(0,1)$, all sufficiently small positive
$\sigma$ have the following properties on this one fixed permutation:
\begin{enumerate}
\item Fresh posterior sampling makes at least $T_k$ mistakes with
      probability at least $1-\eta$.
\item Posterior-label majority makes at least $T_k$ mistakes.
\item The strict $2/3$ confidence criterion fails on at least $T_k$ rounds.
\end{enumerate}
The minimum Boolean circuit size of the target is $\Theta(k)$ for a
fixed complete gate basis of fan-in at most two.
\end{theorem}

The input order depends only on $k$, not on the sampled predictions or on
$\sigma$. The quantifier ``sufficiently small'' may depend on $k,L,\eta$;
we do not give an explicit bound on how small $\sigma$ must be.

\begin{corollary}[A bounded-weight, linear-depth teacher]\label{cor:teacher}
At depth $L=2k+2$, the target in Theorem~\ref{thm:main} has a width-one
implementation with all weights of absolute value at most two and a
strict global output margin greater than $1/2$. Its Gaussian quadratic
energy is less than $(40k+11)/3$.

If $C_L(f)$ is the least positive integer bounding both the teacher's
hidden width and the absolute value of every weight, then $C_L(f)=2$ at
this depth. The learner in the theorem has $N=1$. In particular this is
not a counterexample with $N\ge C_L(f)$.
\end{corollary}

The width qualification is substantive. At $L=2$ the proof still gives
exponentially many mistakes, but the displayed shallow implementation has
exponentially large weights. Increasing the scalar depth gives bounded
weights without changing the function ordering used by the proof. It does
not provide a theorem for a wider learner.

\section{The exact minimum after calibration}

For an affine threshold function $g$, define its full-cube unit-margin cost
\[
 Q(g)=\min\left\{\|A\|_2^2:
       g(\xi) A\mathbin{\cdot}(1,\xi)\ge1
       \text{ for every }\xi\in\{-1,1\}^d\right\}.
\]
The minimum exists and its coefficient vector is unique, by coercivity
and strict convexity. We calculate the minimum over all functions
consistent with a particular history. The calculation first uses only
convex training inequalities; its optimizer then turns out to satisfy the
absolute margin on the entire cube.

Here are the calibration examples, all with selector $z=-1$.
\begin{enumerate}
\item Take diagonal inputs $x=y$ equal to the all-plus vector, each of
its $k$ one-coordinate flips, and the negatives of these vectors.
All these labels are positive.
\item For each $i$, use two negative examples with
$x_i=-1,y_i=+1$, with $x_j=+1,y_j=-1$ for $j<i$, and with the upper equal
bits $x_j=y_j$ all plus in one example and all minus in the other.
Their highest differing bit is $i$, so $X<Y$.
\end{enumerate}
Repeated examples in this description are retained only once. There are
at most $4k+2$ calibration inputs.

For an integer $t\in\{0,\ldots,2^k-1\}$, let $q_t$ be the input with
$z=+1$, $Y=0$ and $X=t$. Its target label is negative.

\begin{lemma}[Unique minimum after a negative query]\label{lem:minimum}
After the calibration and the negative query $q_t$, where $t\ge k$,
the unique minimum-norm affine unit-margin vector fitting these labels is
\[
 A_t(x,y,z)=-t-(t+1)z+\sum_{i=0}^{k-1}2^i(x_i-y_i).
 \tag{1}\label{eq:At}
\]
It also fits every $q_j$ for $j\le t$ and has absolute value at least one
on every cube input. Thus it is also the unique minimum among all
full-cube unit-margin classifiers fitting the calibration and
$q_k,q_{k+1},\ldots,q_t$.
\end{lemma}
\begin{proof}
Write any affine score as
\[
 A(x,y,z)=b+cz+\sum_i a_i x_i+\sum_i h_i y_i,
\]
and put
\[
 u=b-c,\quad v=b+c,\quad
 w_i=\frac{a_i-h_i}{2},\quad s_i=a_i+h_i,\quad U=u-1.
\]
Its squared coefficient norm is
\[
 \|A\|_2^2=\frac{u^2+v^2}{2}
                  +2\sum_i w_i^2+\frac12\sum_i s_i^2.       \tag{2}
\]
The paired positive diagonal examples give
\[
 U\ge0,\qquad
 \left|\sum_j s_j\right|\le U,\qquad
 \left|\sum_j s_j-2s_i\right|\le U.
\]
Therefore $|s_i|\le U$ for every $i$, and
\[
 \sum_i|s_i|\le kU.                                      \tag{3}
\]
Average the two negative carry inequalities at coordinate $i$. The upper
$s_j$ terms cancel, leaving
\[
 u-2w_i+2\sum_{j<i}w_j\le-1.
\]
Induction gives
\[
 w_i\ge\sum_{j<i}w_j+\frac{u+1}{2}
       \ge2^i\left(1+\frac U2\right).                    \tag{4}
\]
At $q_t$, the main-coordinate part of the score is
\[
 2\sum_{i:\,\mathrm{bit}_i(t)=1}w_i
       -\sum_{i:\,\mathrm{bit}_i(t)=0}s_i.
\]
The negative unit-margin constraint, (3), and (4) imply
\[
 v\le-1-2\sum_{\mathrm{bit}_i(t)=1}w_i
             +\sum_{\mathrm{bit}_i(t)=0}s_i
   \le -(2t+1)+(k-t)U
   \le -(2t+1),                                         \tag{5}
\]
where the last inequality uses $t\ge k$. Equations (2)--(5) yield
\[
 \|A\|_2^2\ge
      \frac{1+(2t+1)^2}{2}+2\sum_{i=0}^{k-1}4^i=:Q_t.    \tag{6}
\]
Equality requires $u=1$, $s_i=0$, $w_i=2^i$ and $v=-(2t+1)$,
which is exactly (1).

Conversely, (1) has score $1+2(X-Y)$ on the $z=-1$ branch and
$-(2t+1)+2(X-Y)$ on the $z=+1$ branch. It fits all calibration labels,
and its score at $q_j$ is $2j-2t-1\le-1$ whenever $j\le t$.
Both branch scores are odd integers at every cube input. Its full-cube
absolute margin is therefore at least one, and it attains (6).
\end{proof}

\section{From affine cost to the Bayesian posterior}

\begin{lemma}[Scalar tanh networks preserve the affine cost ordering]
\label{lem:energy}
At every fixed depth $L\ge2$, the minimum Gaussian quadratic energy of a
full-margin function $g$ is a strictly increasing function of $Q(g)$.
After any fixed realizable history whose minimum-$Q$ consistent function
is unique, the strong-prior posterior concentrates on that function.
\end{lemma}
\begin{proof}
The Gaussian quadratic energy is
\[
 \kappa(\theta)=d\|a\|_2^2+\sum_j r_j^2+b^2.
\]
A nonzero scalar tanh chain is strictly monotone and odd, up to an overall
sign. Sign changes of its weights can be absorbed into the first affine
vector without changing the energy. We may therefore take all downstream
weights positive. If the smallest signed first preactivation is $\rho>0$,
the optimal first-layer energy is $dQ(g)\rho^2$.

Let $F(\rho;r,b)$ be the positive scalar output of this chain at first
preactivation $\rho$. The energy infimum over the strict margin event is
\[
 E_L(Q)=\min\left\{dQ\rho^2+\sum_jr_j^2+b^2:
       \rho,r_j,b\ge0,\quad F(\rho;r,b)\ge\tfrac12\right\}. \tag{7}
\]
Coercivity gives a minimizer, and its $\rho$ and downstream weights are
positive because its output is nonzero. If $Q_2>Q_1$, evaluate the
$Q_1$ objective at a minimizer for $Q_2$ to obtain
\[
 E_L(Q_2)\ge E_L(Q_1)+d(Q_2-Q_1)\rho^2>E_L(Q_1).
\]
Strict and weak output margins have the same energy infimum: slightly
increasing the readout turns a weak-margin parameter into a strict-margin
one while approaching its energy.

For each function cell
\[
 U_g=\{\theta:g(x)z_\theta(x)>1/2\text{ at every input}\},
\]
the finite-dimensional Gaussian small-noise upper and lower bounds give
\[
 \lim_{\sigma\downarrow0}\sigma^2\log\Pp_\sigma(U_g)
       =-E_L(Q(g))/2.                                    \tag{8}
\]
The lower bound follows from a ball inside the open cell; the upper bound
uses its closure and coercivity. The matching energy infima just noted
make the two bounds agree. There are finitely many Boolean functions.
Exact label conditioning keeps a function cell whole or deletes it, so
(8) and strict energy separation imply the stated posterior concentration.
The prior's global-margin normalization cancels from these posterior ratios.
\end{proof}

\begin{proof}[Proof of Theorem~\ref{thm:main}]
Present the calibration inputs and then $q_k$. Next present
$q_{k+1},q_{k+2},\ldots,q_{2^k-1}$, and finally all remaining inputs in
any fixed order. At the history ending in $q_t$, Lemma~\ref{lem:minimum}
gives the unique minimum-cost classifier $\sgn A_t$.
Its score at $q_{t+1}$ is $+1$, although the target label is negative.
Lemma~\ref{lem:energy} therefore makes the correct-label posterior
probability tend to zero on every one of these $T_k$ rounds.

For fixed $k,L$, there are only finitely many specified rounds.
Choose a common sufficiently small $\sigma$ so that the correct-label
probability at each is less than both $\eta/T_k$ and $1/3$.
A union bound gives probability at least $1-\eta$ that posterior sampling
is wrong on every specified round. The same probabilities make posterior
majority wrong and the strict $2/3$ criterion fail on all of them.

The target has a linear-size Boolean circuit. In $0/1$ notation,
comparison of the lowest $i+1$ bits obeys the recursion
\[
 G_{-1}=1,\qquad
 G_i=(X_i\wedge\neg Y_i)\ \vee\
           ((X_i\leftrightarrow Y_i)\wedge G_{i-1}).
\]
Each step uses a constant number of bounded-fan-in gates; conjoining
$G_{k-1}$ with the test $z=-1$ uses only a constant further number.
Replacing these gates by NAND gates changes size by a constant factor.
Conversely, all $2k+1$ input bits are essential. For $x_i$, fix $Y=2^i$
and all other $X$ bits to zero; for $y_i$, fix $X=0$ and all other $Y$
bits to zero. Take $z=-1$ in both cases. The selector is essential at
$X=Y$. A fan-in-two circuit with $s$ gates and $d$ essential inputs has
at least $s+d-1$ edges in the connected ancestry of its output and at
most $2s$ edges entering gates. Hence $s\ge d-1=2k$, proving the
matching order lower bound.
\end{proof}

\section{Bounded weights and the width qualification}

\begin{proof}[Proof of Corollary~\ref{cor:teacher}]
Use $A_{2^k-1}$ from (1), divided by $2^k$, as the first preactivation.
Every first-layer coefficient has absolute value at most one, and the
smallest absolute preactivation is $2^{-k}$. Apply the first tanh and
then $2k$ further hidden tanh layers with scalar weight two, followed by
readout weight one. This gives depth $2k+2$.

For $0\le h\le1/2$, concavity and $\tanh(1)>3/4$ give
$\tanh(2h)\ge(3/2)h$. Also $\tanh(2^{-k})\ge(3/4)2^{-k}$.
If the smallest positive activation had not crossed $1/2$ after those
$2k$ amplification layers, it would be at least
\[
 \frac34\,2^{-k}(3/2)^{2k}
     =\frac34(9/8)^k>\frac12,
\]
a contradiction. Once above $1/2$ it stays above $1/2$ under
$h\mapsto\tanh(2h)$. Oddness handles negative scores. This proves the
strict margin and the weight bound two, hence $C_L(f)\le2$. Moreover,
\[
 Q_{2^k-1}=\frac83 4^k-2^{k+1}+\frac13,
\]
so the scaled first vector has squared norm less than $8/3$. Its
Gaussian energy is therefore bounded by
\[
 d\|a\|_2^2+2k\cdot 2^2+1
   <\frac83(2k+1)+8k+1=\frac{40k+11}{3}.
\]
Thus the exponential mistake count occurs despite a linear-energy
neural implementation.

For completeness, $C_L(f)=1$ is impossible. With one neuron in each
layer and all weights bounded by one, the entire downstream scalar chain
is $1$-Lipschitz and vanishes at zero. A strict final margin above $1/2$
therefore forces the first affine margin to be some $\rho>1/2$.
Absorb any overall sign into the first vector and divide it by $\rho$.
The diagonal and carry inequalities in the proof of
Lemma~\ref{lem:minimum} then give normalized
$w_{k-1}\ge2^{k-1}$. In the original first vector,
$w_{k-1}>2^{k-2}>1$ for $k\ge3$, contradicting the bound one on both
corresponding coefficients. Thus $C_L(f)=2$.
\end{proof}

This lower bound concerns a learner that can represent the target with
the same hidden width as the teacher, and with only a constant bound on
the teacher weights. Nevertheless, the condition $N\ge C_L(f)$ used by
some width-uniform learning theorems is stronger and fails here. Adding a
second learner neuron enlarges its function class and can change the
minimum-energy ordering. The present proof gives no reduction from that
wider posterior to the scalar posterior.

\section{Why this does not contradict iid learning bounds}

All Boolean functions represented by a scalar tanh chain, including constants, are affine
threshold functions on the input cube. The number $H_d$ of such Boolean
functions satisfies the hyperplane-arrangement bound
\[
 H_d\le2\sum_{j=0}^{d}\binom{2^d-1}{j},\qquad \log H_d=O(d^2).
\]
For any input distribution, a union bound implies that after
\[
 m\ge\frac{\log H_d+\log(1/\delta)}{\varepsilon}
\]
iid noiseless examples, every consistent classifier has test error at most
$\varepsilon$, with probability at least $1-\delta$.
This controls posterior-sampled test error at every positive prior scale;
constant-factor changes in $\varepsilon$ also control posterior majority
and the strict confidence criterion. No property of the prior masses is
needed for this finite-class argument.

Thus the exhibited learner has polynomial iid sample complexity in input
dimension but exponential total mistakes on the specified order. Its
posterior repeatedly selects a unique cheapest extension of the observed
labels, and the next query reveals where that extension is wrong. The
counterexample relies neither on arbitrary tie-breaking nor on varying
the prior during learning.

\appendix
\section{An all-width consequence for linear activations}
\label{sec:linear}

The width qualification can be removed if the activation is changed to
$\phi(s)=s$. This is a separate result for a deep linear network; it is not
an extension to a wider tanh network. We retain the Boolean domain, the
full-cube output-margin restriction, exact label conditioning, and the
usual layer-scaled independent Gaussian prior.

Let all $L-1$ hidden layers have any common width $N\ge1$. Write
\[
 z_\theta(\xi)=vW^{L-1}\cdots W^1(1,\xi)
              =A_\theta\mathbin{\cdot}(1,\xi),
\]
where $W^1\in\mathbb R^{N\times(d+1)}$, the intermediate matrices are
$N\times N$, and $v\in\mathbb R^{1\times N}$. Independently, the entries
have variances
\[
 \operatorname{Var}(W^1_{ij})=\frac Nd\sigma^2,\qquad
 \operatorname{Var}(W^j_{i\ell})=\sigma^2\quad(2\le j<L),\qquad
 \operatorname{Var}(v_i)=\frac1N\sigma^2.
\]
The quadratic energy is
\[
 \kappa_N(\theta)=\frac dN\|W^1\|_F^2+
       \sum_{j=2}^{L-1}\|W^j\|_F^2+N\|v\|_2^2.
\]
For $N=1$ this agrees with the scalar prior used in the main theorem.

\begin{proposition}[Exponential mistakes at every sufficient width for
linear activations]\label{prop:linear-all-width}
For the target and fixed input permutation in
Theorem~\ref{thm:main}, the same $T_k=2^k-k-1$ forced rounds occur in the
strong-prior limit of this deep linear learner, at every fixed
$L\ge2$ and every fixed width $N\ge1$. All three conclusions of that
theorem hold at every sufficiently small positive scale; the scale
threshold may depend on $N$ as well as $k,L,\eta$.

At depth $L=k+2$, the target has a strict-margin linear-network
implementation with width one and coefficient bound two. Its
bounded-width-and-weight complexity for linear activation is exactly
$C_L^{\mathrm{lin}}(f)=2$. Hence the lower bound holds, in particular,
for every learner width $N\ge C_L^{\mathrm{lin}}(f)$, with no upper bound
on $N$.
\end{proposition}
\begin{proof}
For any desired effective affine vector $A$, the exact minimum energy
of its matrix factorizations is
\begin{equation}\label{eq:linear-factor-energy}
 \min_{\theta:A_\theta=A}\kappa_N(\theta)
       =L\bigl(d\|A\|_2^2\bigr)^{1/L}.
\end{equation}
To prove the lower bound, denote the $L$ factors by $M_1=W^1$,
$M_j=W^j$ for $2\le j<L$, and $M_L=v$, and put
$c_1=d/N$, $c_L=N$, and $c_j=1$ for the intermediate layers. Matrix norm
submultiplicativity and arithmetic--geometric mean give
\[
 \kappa_N=\sum_{j=1}^L c_j\|M_j\|_F^2
 \ge L\left(\prod_{j=1}^L c_j\|M_j\|_F^2\right)^{1/L}
 \ge L\bigl(d\|A\|_2^2\bigr)^{1/L}.
\]
Equality is attained by balanced rank-one factors. For $A\ne0$, choose
one unit vector $e\in\mathbb R^N$ and set
\[
 W^1=s_1e\frac{A}{\|A\|_2},\qquad
 W^j=s_j ee^\top\ (2\le j<L),\qquad v=s_Le^\top,
\]
where
\[
 r=(d\|A\|_2^2)^{1/(2L)},\qquad s_j=r/\sqrt{c_j}.
\]
Their product is $A$, since
$\prod_j s_j=r^L/\sqrt d=\|A\|_2$, and each energy summand is $r^2$.
For $A=0$, zero factors give equality as well.

The least squared norm of an affine vector with signed output margin
at least $1/2$ for a threshold function $g$ is $Q(g)/4$.
Equation~\eqref{eq:linear-factor-energy} therefore yields the exact
function energy infimum
\[
 E_{N,L}^{\mathrm{lin}}(g)
       =L\bigl(dQ(g)/4\bigr)^{1/L}.
\]
This expression is strictly increasing in $Q(g)$ and independent of
$N$. Strict and weak margins have the same infimum by slightly
increasing the final readout. The Gaussian small-noise argument from
Lemma~\ref{lem:energy} applies to these finite-dimensional matrix
parameters: the strict function cell is open, its closure is the
weak-margin cell, and the energy is coercive. Thus the posterior
concentrates on the unique minimum-$Q$ consistent function after each
specified history. Lemma~\ref{lem:minimum} and the same finite-round
union bound prove the claimed mistake statements at every fixed width.
No interchange of the width and small-scale limits is needed.

For the teacher, divide $A_{2^k-1}$ by $2^k$ in the first layer, use
$k$ further hidden scalar layers of weight two, and use readout weight
one. There are $k+1$ hidden layers, so $L=k+2$, and the final output is
exactly $A_{2^k-1}$, with absolute margin at least one. All weights are
bounded by two. Zero-padding to width two proves
$C_L^{\mathrm{lin}}(f)\le2$.

If this complexity were one, the product of the scalar downstream
weights would have absolute value at most one. A strict output margin
above $1/2$ would then force a first affine margin $\rho>1/2$.
Normalize that affine vector by $\rho$ and apply the diagonal and carry
inequalities of Lemma~\ref{lem:minimum}. They force its original
$w_{k-1}$ to exceed $2^{k-2}>1$, contradicting the coefficient bound
one. Hence $C_L^{\mathrm{lin}}(f)=2$.
\end{proof}

This linear-activation result gives a polynomial-iid versus
exponential-online separation even at arbitrarily large learner widths:
the represented Boolean functions are still exactly the affine
threshold functions, so the iid argument in the preceding section is
unchanged. For tanh, additional neurons introduce nonlinear classifiers,
and the factorization identity above no longer applies. The result
with tanh and $N\ge C_L(f)$ therefore remains open.


\paragraph{Verification status.}
The proofs in this note have undergone independent mathematical checks
by collaborating AI agents. They have not been formalized in Lean.
The exact affine calculation is the central step; the probability
statements follow from finite-dimensional Gaussian concentration.
\end{document}
